% 作者题证的编辑转录副本，不是独立下载的上游原件。
% 来源：https://terrytao.wordpress.com/2008/02/04/254a-lecture-9-ergodicity/

\paragraph{Notation.}
Let $(X,\mathcal X,\mu,T)$ be a measure-preserving probability system.
Write $Uf=f\circ T$ for the Koopman operator and
\[
 A_Nf=\frac1N\sum_{n=0}^{N-1}U^nf,
 \qquad \mathcal I=\{E\in\mathcal X:T^{-1}E=E\text{ modulo null sets}\}.
\]
\begin{theorem}[Pointwise ergodic theorem; Tao's Theorem 2]
For every $f\in L^1(X,\mathcal X,\mu)$,
\[
 A_Nf(x)\longrightarrow\E(f\mid\mathcal I)(x)
 \quad\text{for almost every }x.
\]
\end{theorem}

\begin{lemma}[Maximal inequality for the shift; Theorem 1 specialized to $U$]
For real $f\in L^1$, set $Mf=\sup_{N\geq1}A_Nf$. Then
\[
 \lambda\mu\{Mf>\lambda\}\leq\int_{\{Mf>\lambda\}}f\,d\mu
 \qquad(\lambda\in\R).
\]
Consequently, for arbitrary real or complex $f\in L^1$ and $\lambda>0$,
\[
 \mu\left\{\sup_N A_N|f|>\lambda\right\}\leq\frac{\|f\|_1}{\lambda}.
\]
\end{lemma}
\begin{proof}
Since $M(f-\lambda)=Mf-\lambda$, replace $f$ by $f-\lambda$ and reduce to
$\int_{\{Mf>0\}}f\,d\mu\geq0$. For $m\geq1$, define
\[
 F_m=\max_{0\leq N\leq m}\sum_{n=0}^{N-1}U^nf,
\]
where the sum for $N=0$ is zero. Then $F_m\geq0$, $F_m\in L^1$, and
\[
 F_m\leq F_{m+1}=\max(0,f+UF_m).
\]
On $\{F_m>0\}$ the second maximum equals $f+UF_m$, so integration and
positivity of $UF_m$ imply
\[
 \int_XF_m\,d\mu
 \leq\int_{\{F_m>0\}}f\,d\mu+\int_XUF_m\,d\mu.
\]
Measure preservation gives $\int UF_m=\int F_m$, and hence
$\int_{\{F_m>0\}}f\,d\mu\geq0$.
The sets $\{F_m>0\}$ increase to $\{Mf>0\}$. Dominated convergence,
with dominating function $|f|$, proves the first inequality. Apply it to
$|f|$ and use $\int_E|f|\leq\|f\|_1$ to obtain the stated maximal bound.
\end{proof}

\begin{proof}[Proof of the pointwise theorem]
Subtract $\E(f\mid\mathcal I)$, which is invariant, and reduce to proving
\begin{equation}\label{zero-pointwise}
 \limsup_{N\to\infty}|A_Nf(x)|=0\quad\text{almost everywhere}
\end{equation}
when $\E(f\mid\mathcal I)=0$. For a bounded coboundary $f=Ug-g$,
$g\in L^\infty$, the averages telescope:
\[
 A_Nf=\frac{U^Ng-g}{N}.
\]
This proves \eqref{zero-pointwise} for these functions.

The Hilbert-space argument used in the mean ergodic theorem shows that
bounded coboundaries are dense in the subspace of $L^2$ with zero
conditional expectation onto $\mathcal I$. That subspace is itself dense,
in $L^1$ norm, in the subspace of $L^1$ with zero conditional expectation.
Thus the pointwise assertion has been established for an $L^1$-dense class
in the latter subspace.

Choose $f_j$ in this dense class with $\|f-f_j\|_1\to0$. Outside a common
null set, $A_Nf_j(x)\to0$ for all $j$. The triangle inequality gives
\[
 \limsup_{N\to\infty}|A_Nf(x)|
 \leq\sup_N A_N|f-f_j|(x).
\]
By the maximal inequality, for any $\lambda>0$,
\[
 \mu\left\{\limsup_{N\to\infty}|A_Nf|>\lambda\right\}
 \leq\frac{\|f-f_j\|_1}{\lambda}.
\]
The left side does not depend on $j$; taking $j\to\infty$ makes it zero.
Taking positive rational $\lambda$ proves \eqref{zero-pointwise} and hence
the theorem.
\end{proof}
